% GENERATED FILE. Edit canonical lessons, then run npm run generate:books.
% canonical-source-hash: fnv1a64:9230c46a18e0d136
% canonical-lessons: ES-C422-pagina-web, ES-C422-repaso-red, ES-C422-sintesis-buscar

\chapter{The Web Page}
\label{ch:pagina-web}
\hlchaptermodality{422}

\begin{chapteropening}
\textbf{By the end of this chapter:} \emph{I can name a web page, say why it is feminine and why web never takes a plural -s, and read the line a form prints without its article.}

Read the bare noun on a printed form and say the same thing aloud with its article, which is the difference between how a reading paper and a listening paper present the identical word.
\end{chapteropening}

\section[la página web]{la página web — a two-thousand-year-old word doing a thirty-year-old job}

\label{lesson:ES-C422-pagina-web}



\begin{quote}
Say \emph{internet} and \emph{el ordenador}. You have the network and you have the machine. Say what is missing: the thing you actually look at.
\end{quote}

\subsection*{What to know first}
\begin{itemize}
\raggedright
  \item internet — the network itself.
  \item el ordenador — the machine, named in Spanish for what it \emph{orders} rather than what it counts.
  \item el correo electrónico — the other thing that travels over the network.
\end{itemize}

\begin{sounds}
\begin{itemize}
\raggedright
  \item \texttt{stress-antepenultimate} — \emph{PÁ-gi-na}, stress on the third syllable from the end, which is exactly what the written accent is there to tell you. $\to$ reference
  \item \texttt{g-soft} — the \textbf{g} in \emph{página} comes before \textbf{i}, so it is the throaty sound, not the \textbf{g} of \emph{gato}.
\end{itemize}
\end{sounds}

\begin{cousinweb}
\textbf{la página web} — \textquotedblleft{}the web page.\textquotedblright{} Two words from two thousand years apart, and only one of them was borrowed.

\textbf{Página} is Latin \textbf{pagina}, from \textbf{pangere}, \textquotedblleft{}to fasten, to fix in place.\textquotedblright{} A Roman \emph{pagina} was a column of writing, named for the way vine-growers fastened shoots to a trellis: a column of text is writing \textbf{fixed into a frame}.

\noindent
\begin{tabularx}{\linewidth}{@{}>{\raggedright\arraybackslash}X>{\raggedright\arraybackslash}X@{}}
\toprule
\textbf{word} & \textbf{what got fastened} \\
\midrule
\textbf{página} & writing, into a column \\
\textbf{pacto} (a pact) & two parties, into an agreement \\
\textbf{paz} (peace) & the same, after a war \\
\bottomrule
\end{tabularx}

\textbf{Web} is English, taken whole. Spanish did not translate it.
\end{cousinweb}

\begin{grammarlens}[title={Grammar Lens: which half carries the gender}]
\noindent
\begin{tabularx}{\linewidth}{@{}>{\raggedright\arraybackslash}X>{\raggedright\arraybackslash}X>{\raggedright\arraybackslash}X@{}}
\toprule
\textbf{name} & \textbf{article} & \textbf{what the second word is} \\
\midrule
el correo electrónico & \textbf{el} & an adjective, and it agrees with \emph{correo} \\
la página web & \textbf{la} & a noun bolted on, and it never moves \\
\bottomrule
\end{tabularx}

The first word is the real noun and it decides everything. \emph{Web} is a borrowed noun in an adjective's chair, so it never changes: \emph{dos páginas web}, with no \textbf{-s} on \emph{web}.
\end{grammarlens}

\subsection*{Your turn}
Say these aloud:

\begin{itemize}
\raggedright
  \item \textquotedblleft{}la página web\textquotedblright{} — \emph{PÁ-gi-na web}
  \item \textquotedblleft{}la página web está en internet\textquotedblright{}
  \item \textquotedblleft{}dos páginas web\textquotedblright{} — and stop yourself before adding an \textbf{-s} to \emph{web}
  \item which Spanish word for peace comes from the same root as \emph{página} (\emph{paz})
\end{itemize}

\begin{tcolorbox}[breakable,colback=teal!4,colframe=teal!35!black,title={Before you move on}]
What did Latin \emph{pangere} mean? (To fasten, to fix in place.) Why is \emph{página web} feminine? (\emph{Página} is the head noun and it ends in \textbf{-a}.) What happens to \emph{web} in the plural? (Nothing — it stays \emph{web}.)
\end{tcolorbox}
\section[cuatro palabras de internet]{cuatro palabras de internet — four words, four different ways in}

\label{lesson:ES-C422-repaso-red}



\begin{quote}
Name the four things you can now talk about on a screen. Say them with their articles before reading the table.
\end{quote}

\begin{grammarlens}[title={Grammar Lens: borrowed, translated, or built}]
English named all four of these things first. Spanish did something different with each one, and that is the pattern worth carrying away:

\noindent
\begin{tabularx}{\linewidth}{@{}>{\raggedright\arraybackslash}X>{\raggedright\arraybackslash}X>{\raggedright\arraybackslash}X@{}}
\toprule
\textbf{Spanish} & \textbf{English source} & \textbf{what Spanish did} \\
\midrule
internet & internet & took it whole \\
el ordenador & computer & \textbf{refused} it and used a French word instead \\
el correo electrónico & e-mail & translated both halves \\
la página web & web page & translated one half and kept the other \\
\bottomrule
\end{tabularx}

Four things, four strategies, and no rule decided which. \emph{Ordenador} is the interesting one: Spain went to French \emph{ordinateur}, \textquotedblleft{}the thing that puts in order,\textquotedblright{} rather than to the Latin \emph{computare}, \textquotedblleft{}to count.\textquotedblright{} A machine that orders is a different idea about what the machine is for than a machine that counts.
\end{grammarlens}

\begin{tcolorbox}[breakable,skin=enhanced,colback=orange!6,colframe=orange!45!black,boxrule=0.5pt,arc=1mm,left=6pt,right=6pt,top=4pt,bottom=4pt,fonttitle=\bfseries,title={Read this}]
Read aloud. Nothing here is new:

\begin{quote}
El ordenador está en internet. La página web está en el ordenador. El correo electrónico también.
\end{quote}
\end{tcolorbox}

\subsection*{Your turn}
Say these aloud:

\begin{itemize}
\raggedright
  \item the three that take \emph{el} (\emph{el ordenador}, \emph{el correo electrónico}, and the network word, which takes no article at all)
  \item the one that takes \emph{la} (\emph{la página web})
  \item which of the four Spanish refused to borrow (\emph{el ordenador})
  \item \textquotedblleft{}la página web y el correo electrónico\textquotedblright{}
\end{itemize}

\begin{tcolorbox}[breakable,colback=teal!4,colframe=teal!35!black,title={Before you move on}]
Which of the four words is entirely Spanish in both halves? (\emph{El correo electrónico}.) Which came from French rather than English? (\emph{El ordenador}.) Which keeps an English word inside it? (\emph{La página web}.)
\end{tcolorbox}
\section[¿Cuál es la página web?]{¿Cuál es la página web? — the question a form actually asks you}

\label{lesson:ES-C422-sintesis-buscar}



\begin{quote}
Say \emph{la página web} and \emph{el correo electrónico} one after the other. Both are things a form will ask you for.
\end{quote}

\begin{grammarlens}[title={Grammar Lens: the line on a form is the bare noun}]
A printed form drops the article; a person speaking puts it back:

\noindent
\begin{tabularx}{\linewidth}{@{}>{\raggedright\arraybackslash}X>{\raggedright\arraybackslash}X@{}}
\toprule
\textbf{on the form} & \textbf{spoken aloud} \\
\midrule
Correo electrónico: & ¿Cuál es \textbf{el} correo electrónico? \\
Página web: & ¿Cuál es \textbf{la} página web? \\
\bottomrule
\end{tabularx}

This is worth noticing because a reading paper shows you the left-hand column and a listening paper plays you the right-hand one. The same word looks bare in one and dressed in the other.
\end{grammarlens}

\begin{tcolorbox}[breakable,skin=enhanced,colback=orange!6,colframe=orange!45!black,boxrule=0.5pt,arc=1mm,left=6pt,right=6pt,top=4pt,bottom=4pt,fonttitle=\bfseries,title={Read this}]
\begin{quote}
\textbf{Formulario} Nombre: Correo electrónico: Página web:
\end{quote}

Three lines, no articles, and you can now read every one of them.
\end{tcolorbox}

\subsection*{Your turn}
Say these aloud:

\begin{itemize}
\raggedright
  \item \textquotedblleft{}¿Cuál es la página web?\textquotedblright{}
  \item \textquotedblleft{}¿Cuál es el correo electrónico?\textquotedblright{} and hear the article change while the question word does not
  \item \textquotedblleft{}la página web está en el ordenador\textquotedblright{}
\end{itemize}

\begin{tcolorbox}[breakable,colback=teal!4,colframe=teal!35!black,title={Before you move on}]
Why does a form print \emph{Página web:} without \emph{la}? (A form label is the bare noun; the article belongs to speech.) Which article does it take when you say it aloud? (\emph{La}.)
\end{tcolorbox}
